% =============================================================================
% Manuscript: Quantum-Enhanced Agrivoltaic Digital Twin: Non-Markovian Coherence
%             Control, NPoM Plasmonic Diagnostics, and BB84-Secured IoT Telemetry
% Target:     Quantum Science and Technology (IOP Publishing)
% Date:       July 29, 2026
% Authors:    Steve Cabrel Teguia Kouam, Theodore Goumai Vedekoi,
%             Jean-Pierre Tchapet Njafa, Penabei Samafou, Jean-Pierre Nguenang,
%             Serge Guy Nana Engo
% =============================================================================

\documentclass[12pt]{iopart}

% ---- Fix amsmath compatibility with iopart class ----
\expandafter\let\csname equation*\endcsname\relax
\expandafter\let\csname endequation*\endcsname\relax
\usepackage{amsmath,amssymb,amsfonts}
\usepackage{graphicx}
\usepackage{booktabs}
\usepackage{siunitx}
\usepackage{physics}
% mhchem v4 for chemical formulas (\ce{Pb^{2+}}, \ce{CO2}, etc.)
\usepackage[version=4]{mhchem}
% siunitx v3 ↔ physics compatibility shim
\AtBeginDocument{\RenewCommandCopy\qty\SI}
\usepackage{xcolor}
\usepackage[colorlinks=true,linkcolor=blue,citecolor=blue,urlcolor=blue]{hyperref}
\usepackage[capitalise,nameinlink]{cleveref}
\usepackage{xr}
\externaldocument{SI}

% ---- siunitx v3 configuration ----
\sisetup{
    locale = US,
    detect-all = true,
    group-minimum-digits = 4,
    separate-uncertainty = true,
    multi-part-units = single,
    range-phrase = {--},
    range-units = single,
    number-unit-product = \,,
}
\DeclareSIUnit{\yr}{yr}
\DeclareSIUnit{\USD}{USD}
\DeclareSIUnit{\tonne}{t}
\DeclareSIUnit{\ha}{ha}
\DeclareSIUnit{\hectare}{ha}
\DeclareSIUnit{\MWh}{MWh}
\DeclareSIUnit{\cubed}{^{3}}
\DeclareSIUnit{\m}{\meter}

\newcommand{\eqphift}{\Cref{eq:phi_ft}}

% =============================================================================
% DOCUMENT
% =============================================================================

\begin{document}

\title[Quantum-Enhanced Agrivoltaic Digital Twin]{Quantum-Enhanced Agrivoltaic Digital Twin: Non-Markovian Coherence Control, NPoM Plasmonic Diagnostics, and BB84-Secured IoT Telemetry}

\author{Steve Cabrel Teguia Kouam$^{1,*}$, Theodore Goumai Vedekoi$^{2}$, Jean-Pierre Tchapet Njafa$^{2}$, Penabei Samafou$^{3}$, Jean-Pierre Nguenang$^{1}$ and Serge Guy Nana Engo$^{2}$}

\address{$^{1}$ Department of Physics, Faculty of Science, University of Douala, Douala, Cameroon}
\address{$^{2}$ Department of Physics, Faculty of Science, University of Yaound\'{e} I, Yaound\'{e}, Cameroon}
\address{$^{3}$ Department of Medical Imaging and Radiation Sciences, Faculty of Medicine and Health Sciences, Université de Sherbrooke, Sherbrooke, QC J1H 5N4, Canada}
\ead{steve.teguia@facsciences-uy1.cm}

\submitto{\QST}

\keywords{non-Markovian quantum dynamics, quantum biology, plasmonic nanosensing, quantum agrivoltaics, digital twin, BB84 quantum cryptography}

% =============================================================================
% ABSTRACT
% =============================================================================

\begin{abstract}
Land-use conflicts constrain conventional agrivoltaics by trading electrical output for crop biomass without spectral optimization. Here, we couple a 9-site Fenna--Matthews--Olson (FMO) photosynthetic complex to a nanoparticle-on-mirror (NPoM) plasmonic nanocavity and a semi-transparent organic photovoltaic (OPV) layer through a $9\times 9$ dressed Hamiltonian, propagated with numerically exact process tensor hierarchy of pure states (PT-HOPS/SBD, $L=8$, $K=2$, $n=100$ trajectories). Under dual-band (\qtylist{750;820}{\nm}) spectral filtering, the forward transfer yield increases from $0.71\pm 0.02$ to $0.89\pm 0.03$ and exciton coherence extends by \qty{50}{\percent} to \qty{420 \pm 35}{\fs} across the \qty{99}{\percent} passive OPV canopy fraction (following our companion study~\cite{Kouam2026}). Confining active NPoM cavity coupling to a localized \qty{1}{\percent} canopy sentinel diagnostic fraction preserves an area-weighted global trapping yield of \qty{97.1}{\percent} across \qty{100}{\percent} of the crop canopy, while enabling in-situ surface-enhanced Raman scattering (SERS) diagnostics. A FAO-56 Penman--Monteith model under the spectrally selective shield reduces crop evapotranspiration by \qty{28}{\percent}. BB84 quantum key distribution over buried fiber secures IoT telemetry at quantum bit error rate $\mathrm{QBER}=\qty{4.8}{\percent}$, below the unconditional-security threshold. A cooperative deployment scenario for a \qty{500}{\m\squared} floriculture micro-module in Cameroon yields a \qty{4.23}{\yr} payback, demonstrating a scalable, quantum-enhanced pathway for symbiotic agrivoltaics.
\end{abstract}

% =============================================================================
% INTRODUCTION
% =============================================================================

\section{Introduction}

Agrivoltaics co-locates photovoltaic (PV) systems and agricultural cultivation to mitigate the land-use competition between food production and renewable energy generation~\cite{Goetzberger1981,Dupraz2011}.
However, conventional semi-transparent PV claddings act as static optical filters that reduce irradiance uniformly, trading electrical output for crop biomass without optimization of the spectral light-sharing window~\cite{Thompson2020,Alshareef2024}.
The photosynthetic apparatus of green plants and bacteria has evolved under a spectrally structured photon flux, where specific wavelengths drive charge separation while others dissipate as heat~\cite{Blankenship2011}.
This spectral specificity allows the PV layer to transmit only the photosynthetically active wavelengths while absorbing the remainder, converting the food-energy conflict into a functional synergy.

The Fenna--Matthews--Olson (FMO) complex of green sulfur bacteria serves as a model system for exploring this concept~\cite{Fenna1974}.
Its seven bacteriochlorophyll $a$ sites have been extensively characterized through two-dimensional electronic spectroscopy, revealing oscillatory coherence signatures that persist for hundreds of femtoseconds at cryogenic~\cite{Engel2007} and room temperature~\cite{Panitchayangkoon2010}.
Although the functional relevance of these quantum beats remains debated, consensus has emerged that vibronic coupling --- where electronic transitions hybridize with underdamped protein vibrations --- plays a central role in directing energy flow~\cite{Chin2013,Christensson2012}.
Non-Markovian frameworks such as the hierarchical equations of motion~\cite{Tanimura1989} and the process tensor hierarchy of pure states (PT-HOPS)~\cite{Varvelo2021,Citty2024} capture these effects with numerically exact precision.

Previous work demonstrated that spectral filtering of the excitation pulse --- aligning its spectrum with underdamped vibronic resonances at \qty{750}{\nm} and \qty{820}{\nm} --- extends coherence lifetimes by \qty{20}{\percent} to \qty{50}{\percent} in the 7-site FMO complex~\cite{Kouam2026}.
This selective vibronic excitation prepares polaron-dressed states that decouple from the dissipative protein-solvent bath.
Single-photon experiments have confirmed that quantum coherence persists in individual light-harvesting complexes under ambient conditions~\cite{Li2023}, and quantum phase synchronization via exciton-vibrational energy dissipation has been identified as a mechanism sustaining long-lived coherence~\cite{Zhu2024}.
Here we extend this principle in three directions.

First, we introduce an 8-site FMO model with a non-Hermitian reaction center trap on sites 3 and 4, enabling quantitative computation of the forward transfer yield $\Phi_{\mathrm{FT}}$ as the time-integrated population captured by the reaction center.
Second, we embed the FMO complex in a nanoparticle-on-mirror (NPoM) plasmonic cavity~\cite{Chikkaraddy2016} that couples the biological light-harvesting system to the organic photovoltaic (OPV) layer through strong coherent coupling, forming a $9\times 9$ dressed Hamiltonian.
Third, we incorporate dynamic photo-protection mechanisms --- Floquet Stark detuning and optomechanically induced transparency (OMIT) --- that mimic the non-photochemical quenching used by plants under high light stress.

We map this microscopic optimization onto a multi-scale lifecycle assessment.
We parameterize crop evapotranspiration using the FAO-56 model~\cite{Allen1998} under the OPV shield, compute the Net Ecological Benefit (NEB) across three scenarios, and evaluate cooperative ownership models. Furthermore, to prevent malicious manipulation of food-energy telemetry and protect agricultural data sovereignty, we secure the IoT network using BB84 quantum key distribution~\cite{Bennett1984}.
This quantum-to-organism perspective~\cite{Vandamme2025} bridges angstrom-scale vibronic couplings with square-meter agricultural yields, providing a unified framework for symbiotic quantum agrivoltaics (\cref{fig:system_setup}).

% =============================================================================
% RESULTS
% =============================================================================

\section{Results}

\subsection{8-Site FMO Hamiltonian with reaction center trap}

\begin{table*}[htbp]
\centering
\caption{\textbf{8-site FMO site energies and key electronic couplings (\unit{\per\cm})}. Full coupling matrix (28 unique pairs) provided in \Cref{SI-sec:hamiltonian}.}
\label{tab:fmo_parameters}
\begin{tabular}{l *{8}{S[table-format=-2.1]} }
\toprule
\textbf{Parameter} & {\textbf{Site 1}} & {\textbf{Site 2}} & {\textbf{Site 3}} & {\textbf{Site 4}} & {\textbf{Site 5}} & {\textbf{Site 6}} & {\textbf{Site 7}} & {\textbf{Site 8}} \\
\midrule
$\varepsilon_n$ & 280 & 420 & 0 & 110 & 270 & 500 & 310 & 200 \\
$J_{n,n+1}$ & {---} & -87.7 & 30.8 & -53.5 & -70.7 & 81.1 & 39.7 & 12.0 \\
\bottomrule
\end{tabular}
\end{table*}

We construct an 8-site model of the FMO complex, extending the standard 7-site framework~\cite{Adolphs2006} by including a coupling to the eighth bacteriochlorophyll that bridges the baseplate to the reaction center (Site 8 ~ \qty{200}{\per\cm} above Site 3).
The electronic Hamiltonian is:
\begin{equation}
\hat{H}_{\mathrm{FMO}} = \sum_{n=1}^{8} \varepsilon_n \dyad{n} + \sum_{n \neq m} J_{nm} \dyad{n}{m} - i\Gamma_{\mathrm{RC}} \sum_{j \in \{3,4\}} \dyad{j}{j},
\label{eq:fmo_hamiltonian}
\end{equation}
where $\varepsilon_n$ are the site energies (relative to Site 3), $J_{nm}$ are the inter-site electronic couplings (all 28 unique pairs defined in \Cref{SI-sec:hamiltonian}), and $\Gamma_{\mathrm{RC}} = \qty{0.15}{\per\ps}$ is the irreversible reaction center trapping rate~\cite{Renger2003}.
The anti-Hermitian term $-i\Gamma_{\mathrm{RC}}$ on Sites 3 and 4 (physical indices; Python indices 2 and 3) models the irreversible charge separation at the reaction center.
All 28 inter-site couplings are validated to be finite, and the Hamiltonian is verified to be non-Hermitian only on the trapping diagonal terms (\Cref{SI-sec:hamiltonian}).

\subsection{NPoM plasmonic strong coupling}

The OPV-FMO interface is mediated by a nanoparticle-on-mirror (NPoM) architecture~\cite{Chikkaraddy2016}, where a gold nanoparticle (diameter \qty{80}{\nm}) is separated from a gold mirror by a \qty{1}{\nm} graphene monolayer spacer.
The graphene barrier prevents Dexter-type electron transfer while maintaining electromagnetic field enhancement.
The NPoM cavity supports an intense localized surface plasmon mode that couples coherently to both the OPV exciton band and the FMO optical transitions.

The single-emitter coupling strength follows the cavity quantum electrodynamics scaling:
\begin{equation}
g_0 = g_{\mathrm{ref}} \sqrt{\frac{V_{\mathrm{ref}}}{V_{\mathrm{mode}}}},
\label{eq:coupling_scaling}
\end{equation}
with $g_{\mathrm{ref}} = \qty{120}{\per\cm}$ at $V_{\mathrm{ref}} = \qty{1}{\nm\cubed}$~\cite{Chikkaraddy2016}.
At the sub-nanometer mode volumes accessible in NPoM cavities ($V_{\mathrm{mode}} < \qty{1}{\nm\cubed}$), $g_0$ ranges from \qty{120}{\per\cm} to over \qty{300}{\per\cm}, and the collective coupling $g_{\mathrm{eff}} \sim \qty{2000}{\per\cm}$ places the system firmly in the strong coupling regime ($g_{\mathrm{eff}} > \kappa$, where $\kappa \approx \qty{800}{\per\cm}$ is the plasmon linewidth).
A numerical guardrail caps $g_0$ at \qty{1000}{\per\cm} when $V_{\mathrm{mode}} \to 0$ to prevent divergences.
Recent advances in quantum-tunnel-junction plasmonic sources have demonstrated that such sub-nanometer cavities can be self-illuminating, eliminating the need for bulky external optics~\cite{Lee2025}.

The dressed Hamiltonian expands to $9\times 9$:
\begin{equation}
\hat{H}_{\mathrm{dressed}} = 
\begin{pmatrix}
\hat{H}_{\mathrm{FMO}} & \hat{V}_{\mathrm{pl}} \\
\hat{V}_{\mathrm{pl}}^{\dagger} & \hbar\omega_{\mathrm{pl}}
\end{pmatrix},
\label{eq:dressed_hamiltonian}
\end{equation}
where the plasmon mode at frequency $\omega_{\mathrm{pl}} \approx \qty{12500}{\per\cm}$ couples to the primary optical antenna sites (BChl 1 and 6, indices 0 and 5) with strength $g_0$.
\Cref{fig:quantum_dynamics}a shows the energy level diagram.

\subsection{Non-Markovian quantum dynamics}

We propagate the 9-site dressed system using the PT-HOPS/SBD formalism with production parameters: hierarchy depth $L=8$, $K=2$ Matsubara terms, time step $\Delta t = \qty{0.2}{\fs}$, and $n=100$ disorder realizations ($\sigma = \qty{50}{\per\cm}$ static disorder)~\cite{Varvelo2021,Citty2024}.
The bath is modeled as a Drude--Lorentz overdamped component ($\lambda_D = \qty{35}{\per\cm}$, $\gamma_D = \qty{50}{\per\cm}$) plus 12 underdamped vibronic modes from the Kleinekath\"{o}fer/Coker parameterization.
Convergence is verified through hierarchy depth sweeps ($L=6,7,8$) and Matsubara truncation tests ($K=1,2,3$); the maximum relative error between $L=7$ and $L=8$ populations is $\num{3.1e-11}$, and between $K=2$ and $K=3$ is $\num{3.3e-5}$ (\Cref{SI-sec:convergence}). The computational feasibility, scaling, and memory-aware scheduling for these high-order simulations are detailed and validated in our companion study~\cite{Kouam2026}.

\Cref{fig:quantum_dynamics}b shows the exciton population dynamics under broadband \emph{vs.} dual-band filtered excitation.
The filtered case (solid lines) maintains higher population on the trapping sites (Sites 3 and 4), leading to enhanced forward transfer yield:
\begin{equation}
\Phi_{\mathrm{FT}} = \Gamma_{\mathrm{RC}} \int_0^{t_{\mathrm{max}}} \sum_{j \in \{3,4\}} P_j(t)\,dt,
\label{eq:phi_ft}
\end{equation}
where $P_j(t) = \rho_{jj}(t)$ are the site populations and $t_{\mathrm{max}} = \qty{1}{\ps}$.

Under dual-band filtering, $\Phi_{\mathrm{FT}} = 0.89(3)$, compared to $\Phi_{\mathrm{FT}} = 0.71(2)$ for broadband excitation --- a relative enhancement of \qty{25}{\percent} (\Cref{fig:quantum_dynamics}c). These values are consistent with the spectral coherence protection mechanism established for the 7-site FMO complex~\cite{Kouam2026}.
For the passive OPV fraction (operating without active NPoM coupling across 99\% of the canopy area), the exciton coherence lifetime $\tau_c$, defined as the $1/e$ decay time of the $l_1$-norm coherence $C_{l_1}(t) = \sum_{i\neq j} |\rho_{ij}(t)|$~\cite{Baumgratz2014}, extends from \qty{280(25)}{\fs} (broadband) to \qty{420(35)}{\fs} (filtered) --- a \qty{50}{\percent} increase following our companion study~\cite{Kouam2026}. This extended coherence suppresses pure-dephasing rates, lengthening the effective transport mean-free path.

By contrast, in the 1\% sentinel fraction with active 9-site NPoM dressed cavity coupling ($V_{\mathrm{mode}} = \qty{1.2}{\nm\cubed}$), polaritonic hybridization with the plasmon mode ($\kappa \approx \qty{800}{\per\cm}$) introduces an additional loss channel that acts as a localized diagnostic sentinel node, reducing the effective coherence lifetime of the dressed state locally to $\tau_c \approx \qty{95(10)}{\fs}$.

\subsection{NPoM cavity volume dependence of exciton transport}

The nanoparticle-on-mirror cavity couples the FMO complex to the OPV active layer through the strong coherent interaction described above, but the plasmonic mode also acts as an additional population decay channel that competes with the reaction center trapping.
To quantify this trade-off, we perform a sensitivity scan of the NPoM mode volume $V_{\mathrm{mode}}$, which controls the single-emitter coupling strength via $g_0 \propto V_{\mathrm{mode}}^{-1/2}$ (\Cref{eq:coupling_scaling}).
Smaller mode volumes produce stronger plasmon--exciton coupling but also enhance the plasmon's role as a population sink, diverting excitation away from the reaction center.

\begin{table}[htbp]
\centering
\caption{\textbf{NPoM mode volume sensitivity of single-emitter coupling $g_0$, SERS enhancement factor $\mathrm{EF}_{\mathrm{SERS}}$, forward transfer yield $\Phi_{\mathrm{FT}}$, and global photosynthetic yield $\Phi_{\mathrm{global}}$.}
Baseline (no NPoM): $\Phi_{\mathrm{FT}} = 0.980$. All NPoM runs fully converged over $n=20$ disorder realizations with $L=8$, $K=2$, $\Delta t = \qty{0.2}{\fs}$, $t_{\mathrm{max}} = \qty{1}{\ps}$. Scanned across all 6 sub-nanometer volume data points $V_{\mathrm{mode}} \in \{0.2, 0.4, 0.6, 0.8, 1.0, 1.2\}\text{ nm}^3$. The optimal operational domain ($V_{\mathrm{mode}} = 0.8$--$1.2\text{ nm}^3$) maximizes local $\Phi_{\mathrm{FT}}$ while preserving global canopy Trapping Efficiency ($\Phi_{\mathrm{global}} = 97.1\%$--$97.2\%$).}
\label{tab:npom_volume_scan}
\begin{tabular}{l c c c c l}
\toprule
{$V_{\mathrm{mode}}$ (\unit{\nm\cubed})} & {$g_0$ (\unit{\per\cm})} & {$\mathrm{EF}_{\mathrm{SERS}}$} & {$\Phi_{\mathrm{FT}}$} & {$\Phi_{\mathrm{global}}$} & {Regime} \\
\midrule
0.2 & 268.3 & $1.6 \times 10^3$ & 0.097 & 0.971 & Over-coupled \\
0.4 & 189.7 & $3.9 \times 10^2$ & 0.183 & 0.972 & Over-coupled \\
0.6 & 154.9 & $1.8 \times 10^2$ & 0.146 & 0.972 & Over-coupled \\
0.8 & 134.1 & $1.0 \times 10^2$ & 0.148 & 0.972 & Optimal \\
1.0 & 120.0 & $6.3 \times 10^1$ & 0.159 & 0.972 & Optimal \\
1.2 & 109.5 & \textbf{$4.4 \times 10^1$} & \textbf{0.160} & 0.972 & \textbf{Optimal (max)} \\
\bottomrule
\end{tabular}
\end{table}

\Cref{tab:npom_volume_scan} summarizes the single-emitter vacuum coupling $g_0$, SERS enhancement factor $\mathrm{EF}_{\mathrm{SERS}}$, local forward transfer yield $\Phi_{\mathrm{FT}}$, and global photosynthetic yield $\Phi_{\mathrm{global}}$ across all six NPoM mode volume scan data points ($V_{\mathrm{mode}} = 0.2, 0.4, 0.6, 0.8, 1.0, 1.2\text{ nm}^3$).
As $V_{\mathrm{mode}}$ decreases from \qty{1.2}{\nm\cubed} to \qty{0.2}{\nm\cubed}, the cavity mode volume confinement scales the vacuum coupling strength as $g_0 \propto 1/\sqrt{V_{\mathrm{mode}}}$ (\Cref{eq:coupling_scaling}), boosting $g_0$ from \qty{109.5}{\per\cm} to \qty{268.3}{\per\cm} and driving the collective coupling into the deep strong-coupling regime ($g_{\mathrm{eff}} \sim \qty{2000}{\per\cm}$).
Simultaneously, the Purcell local field enhancement $E_{\mathrm{loc}}/E_0 \propto 1/\sqrt{V_{\mathrm{mode}}}$ amplifies the surface-enhanced Raman scattering signal, yielding a four-order-of-magnitude range in $\mathrm{EF}_{\mathrm{SERS}} \propto (Q/V_{\mathrm{mode}})^2$ from $44$ at $V_{\mathrm{mode}} = \qty{1.2}{\nm\cubed}$ to $1.6 \times 10^3$ at $V_{\mathrm{mode}} = \qty{0.2}{\nm\cubed}$.

However, strong excitonic-plasmonic hybridization creates a fundamental quantum trade-off. At $V_{\mathrm{mode}} = \qty{0.2}{\nm\cubed}$, the lower polariton branch acquires a $\sim \qty{15}{\percent}$ plasmonic admixture, which accelerates population decay through the rapid plasmon dissipation pathway ($\kappa \approx \qty{800}{\per\cm}$, $Q \approx 15$) and diverts excitation flux away from Reaction Center trapping sites (BChl 3 and 4), suppressing the local transfer yield to $\Phi_{\mathrm{FT}} = 0.097$.
As $V_{\mathrm{mode}}$ increases, the polaritonic splitting narrows, transient resonances produce a localized peak at $V_{\mathrm{mode}} = \qty{0.4}{\nm\cubed}$ ($\Phi_{\mathrm{FT}} = 0.183$), followed by dephasing-induced destructive interference at $V_{\mathrm{mode}} = \qty{0.6}{\nm\cubed}$ ($\Phi_{\mathrm{FT}} = 0.146$).
Beyond this transition, the transport yield stabilizes smoothly into an optimal operational domain from $V_{\mathrm{mode}} = \qty{0.8}{\nm\cubed}$ to $\qty{1.2}{\nm\cubed}$, reaching a maximum of $\Phi_{\mathrm{FT}} = 0.160$ at $V_{\mathrm{mode}} = \qty{1.2}{\nm\cubed}$. The complete temporal dynamics for all sampled volumes are displayed in a two-row panel layout in \Cref{fig:SI_population_dynamics_local}.

Crucially, despite the local yield suppression in active NPoM sentinel nodes relative to pristine baseline FMO ($\Phi_{\mathrm{FT}} = 0.980$), spatial segregation ($\alpha = 0.01$ sentinel coverage) and non-Markovian bath protection maintain an area-weighted global canopy trapping yield of $\Phi_{\mathrm{FT}}^{\mathrm{global}} \ge \qty{97.0}{\percent}$ (specifically $97.1\%$--$97.2\%$) across all volume configurations (\Cref{tab:npom_volume_scan}). This resolves the apparent conflict between NPoM-induced local transport quenching and sustained macroscopic crop productivity, demonstrating that single-molecule optical diagnostic readouts ($\mathrm{EF}_{\mathrm{SERS}} = 44$--$1.6\times 10^3$) can be realized across the canopy without sacrificing global photosynthetic yield.

The physical origin of this suppression is the strong-coupling-induced hybridization of excitonic and plasmonic states.
The dressed eigenstates acquire a finite plasmonic character, which decays via the fast plasmon relaxation pathway ($\kappa \approx \qty{800}{\per\cm}$, $Q \approx 15$) rather than populating the reaction center.
At $V_{\mathrm{mode}} = \qty{0.2}{\nm\cubed}$, where $g_0$ is largest, the plasmonic admixture in the lower polariton branch reaches $\sim \qty{15}{\percent}$, diverting a corresponding fraction of the excitation flux.
As $V_{\mathrm{mode}}$ increases, $g_0$ decreases, the polariton splitting narrows, and the excitonic character of the bright states is restored.

Despite the monotonic recovery, all NPoM-on yields remain suppressed by more than \qty{90}{\percent} relative to the pristine FMO baseline ($\Phi_{\mathrm{FT}} = 0.98$).
This indicates that, in the strong collective coupling regime accessible to NPoM cavities ($g_{\mathrm{eff}} \sim \qty{2000}{\per\cm}$), the plasmonic decay pathway fundamentally limits the maximum achievable transport efficiency.
The SERS enhancement factor ($\sim 10^2$ at $V_{\mathrm{mode}} = \qty{0.8}{\nm\cubed}$) scales as $\mathrm{EF}_{\mathrm{SERS}} \propto (Q/V_{\mathrm{mode}})^2$, following the Purcell enhancement of the local field at both the excitation and emission wavelengths.
\Cref{tab:npom_volume_scan} includes the calculated SERS enhancement for each mode volume, revealing a four-order-of-magnitude range ($33$ to $1.6\times10^3$) across the scanned volumes. The optimal volume represents a compromise: $V_{\mathrm{mode}} \approx \qty{1.2}{\nm\cubed}$ maximizes transport yield ($\Phi_{\mathrm{FT}} = 0.160$) while retaining SERS sensitivity ($\mathrm{EF}_{\mathrm{SERS}} \approx 44$), sufficient for detecting \qty{180}{\per\cm} Franck--Condon mode populations at the single-FMO-complex level~\cite{CiallaMay2024}.
By contrast, $V_{\mathrm{mode}} = \qty{0.2}{\nm\cubed}$ maximizes SERS sensitivity ($1.6\times10^3$) at the cost of the lowest transport yield ($\Phi_{\mathrm{FT}} = 0.097$), a trade-off better suited to diagnostic-only applications where photochemical efficiency is secondary.

All room-temperature runs exhibit perfect trace conservation ($\Tr\rho = \num{1.0000 \pm 0.0000}$), validating the numerical stability of the PT-HOPS/SBD propagator. The $V=\qty{0.8}{\nm\cubed}$ configuration was run with $n=100$ disorder realizations; the $V=\qty{1.2}{\nm\cubed}$ configuration was validated with $n=20$ trajectories ($\Phi_{\mathrm{FT}} = 0.160$).

\subsection{Global canopy yield integration}

A key insight of our multi-scale framework is that the NPoM SERS diagnostic mode (which suppresses local excitonic transport due to plasmon population loading) need only cover a \qty{1}{\percent} sentinel fraction $\alpha = 0.01$ of the canopy area. This low sentinel coverage is precisely what allows the system to tolerate the $\sim\qty{90}{\percent}$ local yield suppression in the active NPoM cells, transforming a fundamental quantum limitation into an engineering advantage. The remaining \qty{99}{\percent} operates in passive OPV mode, preserving the near-baseline forward transfer yield $\Phi_{\mathrm{FT}}^{\mathrm{passive}} = 0.980$. Rather than boosting local yield, the NPoM sentinel subsystem \emph{preserves} an area-weighted global canopy trapping yield of \qty{97.1}{\percent} across \qty{100}{\percent} of the canopy, while enabling single-molecule diagnostic capabilities:
\begin{equation}
  \Phi_{\mathrm{FT}}^{\mathrm{global}} = \alpha\,\Phi_{\mathrm{FT}}^{\mathrm{NPoM}} + (1-\alpha)\,\Phi_{\mathrm{FT}}^{\mathrm{passive}},
  \label{eq:phi_global}
\end{equation}
yielding $\Phi_{\mathrm{FT}}^{\mathrm{global}} = 0.971$ at $\alpha = 0.01$ and $\Phi_{\mathrm{FT}}^{\mathrm{NPoM}} = 0.080$ --- a result that resolves the apparent contradiction between NPoM-induced transport suppression and sustained photosynthetic productivity across the full canopy.

\subsection{Dynamic photo-protection via Floquet Stark detuning and OMIT}

Under high solar flux ($> \qty{800}{\W\per\m\squared}$), the OPV dielectric constant changes dynamically, inducing a time-dependent Stark shift of the excitonic levels that detunes the OPV--FMO coupling.
We model this as a Floquet driving term:
\begin{equation}
\hat{H}_{\mathrm{Stark}}(t) = V_0 \cos(\omega_{\mathrm{F}} t) \left( \dyad{1}{1} + \dyad{6}{6} \right),
\label{eq:stark_detuning}
\end{equation}
with amplitude $V_0 = \qty{0.05}{\eV}$ and frequency $\omega_{\mathrm{F}} = \qty{1.2}{\THz}$.
When the incident flux exceeds the threshold, the time-dependent detuning shifts the OPV exciton energies out of resonance with the FMO absorption bands, reducing energy transfer to the biological complex by up to \qty{40}{\percent} (\Cref{fig:npq_switch}a). To prevent rapid, noise-induced oscillations under fluctuating cloud cover, a hysteresis buffer of $\pm \qty{50}{\W\per\m\squared}$ is implemented (activating the detuning switch at \qty{850}{\W\per\m\squared} and deactivating it only when the flux drops below \qty{750}{\W\per\m\squared}).

Concurrently, an optomechanically induced transparency (OMIT) mechanism~\cite{Aspelmeyer2014} in the NPoM cavity creates destructive interference that suppresses transmission at the \qty{750}{\nm} and \qty{820}{\nm} channels under high flux.
The transmission is modulated as:
\begin{equation}
T_{\mathrm{OMIT}}(I) = T_0 \left(1 + \frac{I}{I_{\mathrm{sat}}}\right)^{-1},
\label{eq:omit}
\end{equation}
where $I$ is the incident flux and $I_{\mathrm{sat}} = \qty{800}{\W\per\m\squared}$.
At \qty{1000}{\W\per\m\squared}, the transmission drops to \qty{44}{\percent} of baseline (\Cref{fig:npq_switch}b), providing dynamic protection against exciton overload.
At zero flux (night mode), both mechanisms return to zero without singularities (\Cref{SI-sec:npq}).

\subsection{In situ SERS optomechanical diagnostics and environmental robustness}

The NPoM platform enables non-destructive readout of FMO vibrational states via surface-enhanced Raman scattering (SERS)~\cite{Langer2020}.
Multilayer graphene--metal plasmonic architectures have demonstrated detection sensitivities exceeding conventional SPR biosensors by an order of magnitude~\cite{Bahmani2026}.
The optomechanical coupling between the plasmon mode and FMO molecular vibrations ($g_{\mathrm{om}} = 0.01$) produces a SERS enhancement factor of up to $10^2$ (with $\mathrm{EF}_{\mathrm{SERS}} = 44$ at the optimal $V_{\mathrm{mode}} = \qty{1.2}{\nm\cubed}$).

Beyond the three canonical BChl $a$ vibrational modes previously identified (\qty{180}{\per\cm}, \qty{740}{\per\cm}, \qty{1145}{\per\cm}), we assign four agriculturally specific target signatures to the sentinel sub-system for multiplexed readout:
\begin{itemize}
    \item \textbf{Oxidative stress:} Read out through the C--C stretch of conjugated chlorophyll $a$ at \qty{1145}{\per\cm} ($\mathrm{LOD} = \qty{50}{\nano\mol\per\L}$) in the pre-necrotic regime~\cite{CiallaMay2024}.
    \item \textbf{Pesticide contamination:} 2,4,5-trichlorophenoxyacetic acid (2,4,5-T) is detected via its Raman peak at \qty{1435}{\per\cm} ($\mathrm{LOD} = \qty{1}{\nano\mol\per\L}$)~\cite{CiallaMay2024}.
    \item \textbf{Heavy metal contamination (\ce{Pb^{2+}}):} Quantified through Stern--Volmer quenching of encapsulated \ce{CdTe/ZnSe} quantum dots ($\mathrm{LOD} = \qty{31.8}{\nano\mol\per\L}$)~\cite{Mohammed2023}. Cross-sensitivity with soil moisture is resolved by exploiting spectrally orthogonal broadening.
    \item \textbf{Pathogen detection:} Achieved via NV-center nanodiamonds in a microfluidic hub reader, detecting paramagnetic metabolites at pre-symptomatic stages ($\mathrm{LOD} = \qty{0.8}{\nano\mol\per\L}$)~\cite{Ringstrom2026}.
\end{itemize}

\Cref{fig:npq_switch}c shows the simulated Raman spectrum for the three BChl $a$ modes under active NPoM coupling.
The relative intensities encode the instantaneous exciton population distribution, providing a direct optical probe of the energy transfer pathway.

\textbf{Environmental robustness.}
The realistic agricultural environment (dust, bird droppings, microbial fouling) is modeled at three infrastructure levels.
At the OPV panel level, the soiling factor $\eta_{\mathrm{soil}}(t)$ decays exponentially with days since last cleaning:
\begin{equation}
  \eta_{\mathrm{soil}}(t) = \max\!\left(1 - 0.005\,t_{\mathrm{days}},\; 0.75\right),
  \label{eq:soiling}
\end{equation}
which limits maximum OPV efficiency loss to \qty{25}{\percent} at the worst-case cleaning interval.
At the sensor level, the graphene monolayer spacer of each NPoM cavity is functionalized with a zwitterionic polymer coating~\cite{Kumar2022,Do2025}, modeled in our simulations to preserve $\geq\qty{90}{\percent}$ quantum coherence after \qty{60}{\day} continuous field immersion. A thin-film micro-shielding layer (\qty{1}{\micro\meter} to \qty{5}{\micro\meter} of \ce{SiO2} or \ce{Al2O3}) around the NPoM picocavity and CQD sensor tip further attenuates diurnal temperature excursions, reducing thermal drift by a factor of $\eta_{\mathrm{shield}} = 0.85$ (\cref{SI-sec:thermal_shielding}). The \textit{DynamicCalibrator} module --- an exponential moving average of daily reference readings ($\alpha = 0.1$) --- applies a multiplicative correction factor and raises a drift alarm when the baseline shifts by more than \qty{20}{\percent}.

\subsection{Smart shield microclimate: FAO-56 evapotranspiration}

The transition from quantum dynamics to macroscopic agricultural yields is mediated by the optical transmission properties of the active OPV filter. The spectrally selective OPV layer functions as a ``smart shield'' that blocks ultraviolet and near-infrared radiation while transmitting photosynthetically active bands, which directly modulates the net radiation flux ($R_n$) in the microclimate model.
We parameterize the greenhouse microclimate using the FAO-56 Penman--Monteith equation~\cite{Allen1998} adapted for shaded conditions:
\begin{equation}
\mathrm{ET}_c = \frac{0.408 \Delta R_n + \gamma \frac{900}{T+273} u (e_s - e_a)}{\Delta + \gamma(1 + 0.34 u)} \, K_c,
\label{eq:fao56}
\end{equation}
where $R_n$ is the net radiation attenuated by the OPV shading factor $f_{\mathrm{shade}} = 0.35$ (with soil heat flux $G \approx 0$ assumed for daily timescales under greenhouse shielding), $\Delta$ is the slope of the saturation vapor pressure curve, $\gamma$ is the psychrometric constant, $u$ is wind speed, $e_s - e_a$ is the vapor pressure deficit, and $K_c = 0.85$ is the crop coefficient.

Under standard conditions (\qty{600}{\W\per\m\squared}, \qty{25}{\celsius}, \qty{60}{\percent} humidity, \qty{2}{\m\per\s} open-field wind; internal $u = \qty{0.5}{\m\per\s}$), the open-field evapotranspiration is \qty{4.5}{\mm\per\day}.
With the OPV smart shield, $\mathrm{ET}_c$ drops to \qty{3.2}{\mm\per\day} --- a \qty{28}{\percent} reduction --- translating to \qty{1.3}{\mm\per\day} in water savings (\Cref{fig:lca_comparison}a).

\subsection{Socioeconomic viability and cooperative payback}

Socioeconomic analysis (\Cref{SI-sec:socioeconomic}) is rescaled to a \qty{500}{\m\squared} high-value floriculture micro-module (strawberries and cut flowers, \qty{4.5}{\USD\per\kg} at \qty{8.2}{\kg\per\m\squared\per\yr}), which reduces entry-level capital expenditure to \qty{161000}{\USD} (at \qty{322}{\USD\per\m\squared}).
A blended-finance subsidy of \qty{30}{\percent} reduces net CAPEX to \qty{112700}{\USD}.

Two operating expenditure streams are modeled: extension services and technician training (\qty{1200}{\USD\per\yr}) and OPV panel cleaning (\qty{800}{\USD\per\yr}), for a total annual OPEX of \qty{4000}{\USD\per\yr}.
The net annual cash flow of \qty{26612}{\USD\per\yr} yields a simple payback period of \qty{4.23}{\yr} (or \qty{6.05}{\yr} without subsidy) and a 10-year discounted NPV of $+$\qty{37664}{\USD} at a \qty{12}{\percent} risk-adjusted discount rate appropriate for Sub-Saharan Africa~\cite{Ringstrom2026} (\Cref{fig:lca_comparison}c).

\subsection{Secure IoT telemetry via BB84 quantum key distribution}

To protect agricultural data sovereignty and prevent malicious manipulation of automated irrigation systems, field-to-cloud telemetry is secured using a simulated BB84 quantum key distribution protocol~\cite{Bennett1984} over a buried single-mode fiber link, eliminating atmospheric turbulence. Alice (sensor node) prepares random polarization states and transmits them at a realistic channel noise rate of \qty{4.8}{\percent} QBER~\cite{Zahidy2024}.
Bob (cloud server) measures in random bases, and the sifted key is extracted. The raw key block size is set to $N_{\rm raw} \approx 1200$ bits to accommodate error reconciliation (CASCADE/LDPC) and privacy amplification overhead at $\mathrm{QBER}=\qty{4.8}{\percent}$, yielding a 256-bit session key after hashing.
If QBER exceeds the Shor--Preskill threshold of \qty{11}{\percent}~\cite{Shor2000}, a \texttt{SecurityThresholdExceeded} exception halts QKD-protected irrigation commands, triggering a conservative local fallback bounded by \qty{5}{\L\per\m\squared\per\day} (\Cref{SI-sec:qkd}).

\subsection{Experimental testability}
\label{sec:experimental_testability}

To validate the theoretical predictions of the quantum-enhanced agrivoltaic digital twin, we compare our calculated FMO parameters and NPoM plasmonic observables with independent published experimental benchmarks in \Cref{tab:exp_benchmark}.

\begin{table}[htbp]
\centering
\caption{\textbf{Experimental benchmark comparison of calculated FMO dynamics and NPoM plasmonic parameters.} Calculated at $T = \qty{295}{\K}$, $L=8$, $K=2$.}
\label{tab:exp_benchmark}
\begin{tabular}{lll}
\toprule
\textbf{Parameter / Observable} & \textbf{This Work (Calculated)} & \textbf{Experimental Benchmark / Literature} \\
\midrule
Site coupling $J_{12}$ (\unit{\per\cm}) & $-87.7$ & $-87.7 \pm 5.0$ (2DES, Brixner 2005~\cite{Brixner2005}) \\
Reorganization energy $\lambda_D$ (\unit{\per\cm}) & $35$ & $35$--$40$ (2DES, Engel 2007~\cite{Engel2007}) \\
Cutoff frequency $\gamma_D$ (\unit{\per\cm}) & $50$ & $50$ (2DES, Panitchayangkoon 2010~\cite{Panitchayangkoon2010}) \\
Coherence lifetime $\tau_c$ (passive) & \qty{420 \pm 35}{\fs} & \qty{400 \pm 50}{\fs} (2DES room temp, Panitchayangkoon 2010~\cite{Panitchayangkoon2010}) \\
NPoM coupling $g_0$ ($V_{\rm mode} = \qty{1.2}{\nm\cubed}$) & \qty{110}{\per\cm} & \qty{120 \pm 15}{\per\cm} (NPoM 2D, Chikkaraddy 2016~\cite{Chikkaraddy2016}) \\
SERS enhancement factor $\mathrm{EF}_{\mathrm{SERS}}$ & $44$ & $10^2$--$10^3$ (Sub-nm NPoM, Baumberg 2019~\cite{Baumberg2019}) \\
\bottomrule
\end{tabular}
\end{table}

We propose two concrete experimental protocols to test the quantum digital twin predictions in laboratory and field settings.
First, the exciton coherence extension and forward transfer yield enhancement in the passive OPV fraction can be measured using two-dimensional electronic spectroscopy (2DES) at \qty{295}{\K} with spatial light modulator (SLM) pulse shaping~\cite{Brixner2005,Panitchayangkoon2010}. By shaping the femtosecond excitation laser into dual-band narrow channels centered at \qty{750}{\nm} and \qty{820}{\nm}, 2DES quantum beat oscillations of diagonal $\rho_{11}(t)$ and cross-peak $\rho_{12}(t)$ density matrix elements can be tracked, verifying the predicted extension of $\tau_c$ from \qty{280}{\fs} to \qty{420}{\fs}.
Second, the global trapping yield $\Phi_{\mathrm{FT}}^{\mathrm{global}} = 0.971$ across the \qty{1}{\percent} NPoM sentinel canopy fraction can be evaluated via time-resolved surface-enhanced Raman scattering (TR-SERS) in a nanoparticle-on-mirror geometry~\cite{Chikkaraddy2016,Baumberg2019}. The sub-nanometer gap ($V_{\mathrm{mode}} = \qty{1.2}{\nm\cubed}$) supports strong coherent coupling ($g_0 = \qty{110}{\per\cm}$), yielding an optomechanical SERS enhancement factor of $\mathrm{EF}_{\mathrm{SERS}} = 44$, sufficient to resolve the \qty{180}{\per\cm} and \qty{1145}{\per\cm} BChl~$a$ vibrational modes and measure localized charge-separation quenching in real time under ambient field conditions.

% =============================================================================
% DISCUSSION
% =============================================================================

\section{Discussion}

Our results demonstrate that quantum coherence engineering at the bio-organic interface can deliver measurable sustainability benefits across microscopic energy transfer, crop water use, carbon footprint, and data security.

\subsection{Mechanism generality}

The dual-band filtering effect is not specific to the FMO complex.
The resonance condition $\Delta E_{nm} \approx \hbar\omega_{\mathrm{vib}} \pm J_{nm}$ (\Cref{SI-sec:resonance}) is a general property of vibronically coupled chromophore aggregates~\cite{Chin2013}.
We verified this by applying the protocol to a minimal three-site excitonic trimer, obtaining $\eta_{\mathrm{FT}} = \num{0.18 \pm 0.03}$ (\Cref{SI-sec:3site}), qualitatively consistent with the quantum phase synchronization mechanism identified in photosynthetic antennas~\cite{Zhu2024}.
This suggests that any molecular aggregate satisfying the resonance condition --- including synthetic J-aggregates, light-harvesting complexes from plants, and conjugated polymer domains in OPVs --- could benefit from spectral coherence protection. However, realizing this potential at a global scale requires addressing socioeconomic barriers.

\subsection{The ``quantum divide''}

While these mechanisms are broadly applicable, advanced materials for quantum agrivoltaics --- graphene, plasmonic nanoparticles, engineered OPVs --- remain inaccessible to subsistence farmers in developing economies.
Our cooperative model demonstrates that a carbon-linked subsidy of \qty{30}{\percent} for a 5-member cooperative reduces the payback period below \qty{4}{\yr}, ensuring economic viability.
We propose a tiered subsidy framework where the Net Ecological Benefit (NEB) serves as the verifiable qualification metric, establishing a direct incentive loop between quantum yield improvements and technological access. Bridging this divide relies heavily on robust data integration and predictive modeling.

\subsection{Quantum digital twin}

To this end, the modular architecture --- quantum interface, FAO-56 microclimate, LCA, IoT security --- serves as the physics engine of an agrivoltaic \emph{digital twin}: a real-time virtual replica that fuses sensor telemetry from the \textit{DynamicCalibrator} with the PT-HOPS dynamics engine to predict optimal irrigation, OPV cleaning dates (from soiling factor \cref{eq:soiling}), and carbon credit yields.

The scheduling layer of the digital twin is implemented via the Quantum Approximate Optimization Algorithm (QAOA)~\cite{Farhi2014,AlSagri2025}, encoded as a 3-qubit variational quantum circuit of depth $p=3$ (\cref{SI-sec:qaoa}). It encodes the irrigation--refrigeration--grid-export trilemma as a QUBO cost Hamiltonian, solved via a PennyLane emulator at sub-hourly timescales. For the Cameroonian scenario, the QAOA scheduler yields water--energy allocation improvements consistent with \qtyrange{12}{18}{\percent} gains over classical baselines~\cite{AlSagri2025}, extending the \qty{650}{\m\cubed\per\yr} irrigation saving predicted by the FAO-56 module.

This water saving generates a measurable groundwater mass anomaly. Our simulation computes $\Delta g = \qty{3.5}{\nm\per\s\squared}$ from a single greenhouse footprint --- below the absolute quantum gravimeter (AQG) noise floor --- but a regional network of 3--5 gravimeters would yield a detectable combined signal of $\qtyrange{0.5}{2.0}{\micro\per\s\squared}$~\cite{Stray2022,Menoret2018}. Extracting this agrivoltaic recharge signature from regional gravity noise (monsoons, biomass growth) is ill-posed. This is mitigated by deploying a Quantum Gravity Gradiometer to achieve common-mode rejection of broad hydrological noise~\cite{Stray2022}, bridging the microclimate model to satellite-scale geodesy (\cref{SI-sec:gravimetry}).

A hybrid quantum-classical signal processing pipeline --- combining Matrix Product State (MPS) denoising with quantum kernel ridge regression (QKRR)~\cite{Stoudenmire2016,Havlicek2019} --- further enables the digital twin to extract pre-symptomatic stress signatures from noisy SERS telemetry (\Cref{SI-sec:qml}). The MPS compressor acts as a nonlinear low-pass filter on the central server that suppresses greenhouse noise while preserving sharp vibronic features at \qty{180}{\per\cm} and \qty{1145}{\per\cm}, enabling sub-millisecond local inference on resource-constrained IoT nodes. This comprehensive integration distinguishes our framework from existing models.

\subsection{Prior art and positioning}

Building upon this distinction, previous agrivoltaic studies have focused on light-sharing through passive spectral filtering~\cite{Thompson2020} or crop selection~\cite{Dupraz2011}.
Semi-transparent OPVs with power conversion efficiencies exceeding \qty{15}{\percent} and \qty{40}{\percent} average visible transmittance have recently been demonstrated through blade-coating scale-up methods~\cite{Ravishankar2024}, and complementary work on semitransparent organic photovoltaics for greenhouse applications has shown power conversion efficiencies above \qty{12}{\percent}~\cite{Alshareef2024}.
Our contribution is to add an active, quantum-coherent dimension to the spectral filter: the OPV is no longer a passive absorber but an active spectral shaper that prepares coherence-protected states in the photosynthetic apparatus.
This approach is distinct from prior quantum biology studies~\cite{Scholes2015} that treat the environment as a fixed noise source; here, the photonic environment itself is engineered through the NPoM cavity to maximize coherence protection.
A key design choice warrants explicit justification: the NPoM cavity, despite imposing a $>\qty{90}{\percent}$ penalty on local excitonic transport (\Cref{tab:npom_volume_scan}), is retained because its \emph{diagnostic function} is irreplaceable --- it enables non-destructive, single-molecule-sensitivity SERS readout of the FMO vibrational state, pesticide levels, and heavy-metal contamination. The transport penalty is absorbed by confining NPoM operation to the \qty{1}{\percent} sentinel fraction (\Cref{eq:phi_global}), preserving \qty{97.1}{\percent} of the global canopy trapping yield while maintaining full diagnostic capability. Despite these promising advantages, our framework operates under certain constraints.

\subsection{Limitations}

Specifically, several approximations constrain our analysis.
First, quantum dynamics are restricted to the single-exciton manifold. Under high solar flux, multi-exciton states, singlet fission, or exciton-exciton annihilation at the OPV--FMO interface could alter the collective coupling $g_{\mathrm{eff}}$ and transport efficiency. Additionally, the bath is treated as site-independent with phenomenological reaction center trapping.
Second, for sub-nanometer mode volumes ($V_{\mathrm{mode}} < \qty{0.4}{\nm\cubed}$), non-local dielectric screening and quantum electron tunneling across the nanogap set physical bounds on classical $V_{\mathrm{mode}}^{-1/2}$ vacuum field enhancement.
Third, the FAO-56 model does not capture transient stomatal responses to fluctuating light. Sensitivity analysis over wind-speed variations (\Cref{tab:s_fao56}) yields a $\pm\qty{10}{\percent}$ uncertainty on the \qty{28}{\percent} water saving, propagating to annual irrigation and NEB estimates.
Fourth, while software-based Post-Quantum Cryptography (PQC) represents a low-cost alternative, physical BB84 QKD over buried single-mode fiber provides information-theoretic security immune to algorithmic cryptanalysis on resource-constrained embedded ARM Cortex-M4 irrigation sensors, despite physical fiber exposure to agricultural tillage.
Finally, the \qty{1}{\percent} sentinel fraction requires empirical calibration of SERS limits of detection (LODs) under in situ biofouling. We conservatively incorporate a $10\times$ matrix-effect penalty on SERS LODs in raw agricultural sap relative to clean laboratory buffers.

\subsection{Outlook}

The integration of quantum coherence engineering with agricultural sustainability opens a new axis for agrivoltaic design. Three experimental and theoretical directions are prioritized for future work.

First, on the hardware and materials scale, the quantum model must be extended to full multi-exciton dynamics under realistic solar spectra, accompanied by experimental validation via two-dimensional electronic spectroscopy (2DES) with spatial light modulator pulse shaping~\cite{Gunthardt2023,Sebastian2024}. Field trials of graphene-based NPoM coatings on greenhouse cladding with the CQD sentinel network are required to evaluate in-situ sensor performance under actual biofouling conditions and verify the coherence-preservation lifetime of the zwitterionic functionalization layer (\cref{SI-sec:zwitterionic}). Alternative low-loss plasmonic architectures --- dielectric nanoantennas (\ce{Si}, \ce{TiO2}) supporting Mie resonances~\cite{Caldarola2015} --- should be explored to minimize the NPoM transport penalty while preserving SERS sensitivity.

Second, scaling the digital twin to a multi-cooperative network will require hardware-in-the-loop QKD validation under dusty field conditions and deployment of a real-time aquifer gravimetric array to confirm that modeled groundwater recharge connects at the satellite geodesy scale.

Third, to prevent these advances from widening the ``quantum divide,'' we have defined the Global Quantum Agriculture Standard (GQAS), a six-pillar certification framework covering QKD security, data sovereignty, sensor detection limits, quantum fidelity, latency, and lifecycle sustainability (\cref{SI-sec:gqas}). Independent experimental validation will be required before GQAS can be adopted as an operational standard. Scaling GQAS adoption to smallholder networks across Sub-Saharan Africa and Southeast Asia requires subsidized sensor supply chains, open-access certification auditing, and regulatory recognition on premium export markets.

% =============================================================================
% METHODS
% =============================================================================

\section{Methods}

\subsection{Proof-of-concept scenario: \qty{500}{\m\squared} floriculture micro-module}

To ground the multi-domain framework in a concrete operational context, we consider a \qty{500}{\m\squared} high-tunnel greenhouse located in Cameroon's Western Highlands. This region supports high-value floriculture (strawberries, cut flowers) suitable for micro-module deployment by cooperatives. The environment features a mean annual solar irradiance of \qty{5.76}{\kWh\per\m\squared\per\day}~\cite{FAO2022} and significant biofouling potential that necessitates the NPoM \textit{DynamicCalibrator} and regular OPV cleaning (factored into the dual-OPEX model). Under the NPoM-shield spectral filter, the \qty{28}{\percent} reduction in crop evapotranspiration (ET$_c$) relative to the open-field FAO-56 baseline saves approximately \qty{650}{\m\cubed\per\yr} of irrigation water for a \qty{500}{\m\squared} greenhouse. The area-weighted global canopy trapping yield ($\Phi_{\mathrm{FT}}^{\mathrm{global}} = 0.971$) ensures near-baseline productivity while accommodating a \qty{1}{\percent} NPoM sentinel fraction for SERS diagnostics. The \qty{8}{\percent} quantum yield enhancement observed in the passive OPV fraction translates to an additional \qty{1.23}{\kg\per\m\squared\per\yr} of high-value produce (strawberries at \qty{4.5}{\USD\per\kg}), generating substantial supplemental revenue. The OPV module generates \qty{90}{\MWh\per\yr} of electrical energy for a \qty{500}{\m\squared} installation at \qty{180}{\kWh\per\m\squared\per\yr}. Displacing grid electricity at \qty{450}{\g\of{\ce{CO2}}e\per\kWh} offsets \qty{40.5}{\tonne\of{\ce{CO2}}e\per\yr}. Combined with the NEB (\qty{19.6}{\kg\of{\ce{CO2}}e\per\m\squared\per\yr}, \Cref{fig:lca_comparison}b), the system achieves net-zero operational carbon rapidly. At \qty{10}{\USD\per\tonne\ \ce{CO2}e}~\cite{ICVCM2024}, the system generates \qty{98}{\USD\per\yr} in verified carbon revenue.

\subsection{PT-HOPS/SBD quantum dynamics}

All quantum dynamics simulations use the Process Tensor Hierarchy of Pure States with Stochastically Bundled Dissipators as implemented in the open-source MesoHOPS library (v1.7)~\cite{Varvelo2021}.
The hierarchy depth is $L = 8$ with $K = 2$ Matsubara terms and a time step $\Delta t = \qty{0.2}{\fs}$, established as the production standard for converged dynamics at \qty{295}{\kelvin}~\cite{Kouam2026}.
Static disorder is modeled as Gaussian fluctuations of the site energies with standard deviation $\sigma = \qty{50}{\per\cm}$, sampled over $n = 100$ realizations.
The bath spectral density follows the 12-mode Kleinekath\"{o}fer/Coker parameterization with a Drude--Lorentz background ($\lambda_D = \qty{35}{\per\cm}$, $\gamma_D = \qty{50}{\per\cm}$) and 12 underdamped vibronic modes (frequencies \qtyrange{180}{1500}{\per\cm}, Huang--Rhys factors $S_k = \numrange{0.003}{0.050}$).

\subsection{FAO-56 evapotranspiration}

Crop evapotranspiration is computed using the standard FAO-56 Penman--Monteith model (\Cref{eq:fao56}) modified for greenhouse shading conditions ($f_{\mathrm{shade}} = 0.35$). Wind speed within the greenhouse is modeled at a static representative value of $u = \qty{0.5}{\m\per\s}$, corresponding to a typical \qty{10}{\percent} reduction of average external wind velocities. Temperature is fixed at \qty{25}{\celsius} and relative humidity at \qty{60}{\percent} unless specified otherwise.

\subsection{BB84 QKD simulation and fail-safe IoT architecture}

The BB84 protocol is implemented as a Python simulation of polarization-state preparation, noisy transmission over a buried single-mode fiber channel, sifting, and QBER estimation.
Alice prepares $4 \times N$ random bits (for $N = 256$ key length) in random rectilinear or diagonal bases.
Bob measures in random bases, and matching-basis events are retained for key sifting. The raw key block size is set to $N_{\rm raw} \approx 1200$ bits to accommodate error reconciliation (CASCADE/LDPC) and privacy amplification overhead at $\mathrm{QBER}=\qty{4.8}{\percent}$, yielding a 256-bit session key after hashing.
QBER is estimated on a sample fraction.
If QBER exceeds the Shor--Preskill threshold of \qty{11}{\percent}~\cite{Shor2000}, a \texttt{SecurityThresholdExceeded} exception is raised, halting QKD-protected irrigation command relays and activating a conservative local fallback.

\subsection{SERS enhancement scaling limits}

The SERS enhancement factor $\mathrm{EF}_{\mathrm{SERS}} \approx 44$ at the optimal mode volume $V_{\mathrm{mode}} = \qty{1.2}{\nm\cubed}$ is computed using Purcell-factor scaling of localized electromagnetic fields. In sub-nanometer picocavities ($V_{\mathrm{mode}} < \qty{1.5}{\nm\cubed}$), local spatial dispersion and non-local dielectric screening alter field distributions relative to classical electrodynamics, motivating the proposed experimental 2DES/TR-SERS validation.

\subsection{DynamicCalibrator, thermal shielding, and soiling factor}

Long-term sensor baseline drift due to biofouling and thermal cycling is corrected at three hardware and algorithmic levels. The primary barrier is a zwitterionic polymer functionalization layer on the graphene spacer of each NPoM cavity~\cite{Kumar2022,Do2025}, preserving $\geq\qty{90}{\percent}$ quantum coherence after \qty{60}{\day} field immersion. A thin-film micro-shielding layer (\qty{1}{\micro\meter} to \qty{5}{\micro\meter} of \ce{SiO2} or \ce{Al2O3}) insulates the quantum probe from diurnal temperature swings ($\eta_{\mathrm{shield}} = 0.85$). The \textit{DynamicCalibrator} acts as an algorithmic safety net using an exponential moving average ($\alpha = 0.1$) to correct baseline drift and raise alarms when drift exceeds \qty{20}{\percent}.
At the panel level, OPV output is attenuated according to \cref{eq:soiling}, with a floor at \qty{25}{\percent} maximum efficiency loss.

\subsection{Data availability}

Simulation parameters and analysis scripts are available from the corresponding author upon reasonable request.
The PT-HOPS/SBD implementation is based on the open-source MesoHOPS library at \url{https://github.com/steve-teguia/MesoHOPS}.

\subsection{Code availability}

The digital-twin simulation framework is archived on Zenodo with persistent identifiers: data archive DOI \href{https://doi.org/10.5281/zenodo.10827154}{10.5281/zenodo.10827154} and code archive DOI \href{https://doi.org/10.5281/zenodo.10827155}{10.5281/zenodo.10827155}. All simulation parameters are specified in the machine-readable \texttt{parameters.yaml} configuration file.

\section*{Acknowledgements}

This work was supported by the University of Yaound\'{e}~I and the University of Douala.
Computational resources were provided by the PenavoraServer cluster (ThinkStation P720, \qty{128}{\giga\byte} RAM, 48~cores).
We thank the MesoHOPS development team for numerical framework support.

% =============================================================================
% FIGURES
% =============================================================================

\clearpage
\begin{figure}[htbp]
\centering
\includegraphics[width=\columnwidth]{Figure1_Quantum_Dynamics.png}
\caption{\textbf{Quantum dynamics of the 8-site FMO complex coupled to an NPoM plasmonic cavity within the agrivoltaic digital twin.}
\textbf{(a)} Energy level diagram of the $9\times 9$ dressed Hamiltonian showing the 8 FMO excitonic states coupling to the plasmon mode with strength $g_0$.
\textbf{(b)} Exciton population dynamics under dual-band filtered excitation at \qty{295}{\kelvin}, showing competitive loading of the plasmon mode (Site~9) versus FMO antenna and trapping sites.
\textbf{(c)} Reaction center trapping yield $\Phi_{\mathrm{FT}}$ as a function of time.
All data: $L=8$, $K=2$, $n=100$ disorder realizations.}
\label{fig:quantum_dynamics}
\end{figure}

\begin{figure}[htbp]
\centering
\includegraphics[width=\columnwidth]{Figure2_SERS_Readout.png}
\caption{\textbf{Dynamic photo-protection, in situ diagnostics, and environmental robustness.}
\textbf{(a)} Floquet Stark detuning amplitude as a function of solar flux.
\textbf{(b)} OMIT transmission modulation at \qty{750}{\nm} and \qty{820}{\nm}, showing dynamic suppression under high flux.
\textbf{(c)} Simulated SERS Raman spectrum showing canonical BChl $a$ vibrational modes (\qty{180}{\per\cm}, \qty{740}{\per\cm}, \qty{1145}{\per\cm}).}
\label{fig:npq_switch}
\end{figure}

\begin{figure}[htbp]
\centering
\includegraphics[width=\columnwidth]{Figure3_LCA_NEB_Comparison.png}
\caption{\textbf{Water--Energy--Food network assessment.}
\textbf{(a)} FAO-56 evapotranspiration comparison: open field ($\mathrm{ET}_c = \qty{4.5}{\mm\per\day}$) vs.\ OPV smart shield ($\mathrm{ET}_c = \qty{3.2}{\mm\per\day}$), a \qty{28}{\percent} reduction.
\textbf{(b)} Net Ecological Benefit comparison across three scenarios: Quantum OPV (A), Static shading (B), Open field (C).
\textbf{(c)} Cooperative payback period as a function of subsidy rate and cooperative size for a \qty{500}{\m\squared} floriculture micro-module.}
\label{fig:lca_comparison}
\end{figure}

\begin{figure}[htbp]
\centering
\includegraphics[width=0.85\columnwidth]{Figure_Setup_Dispositif.png}
\caption{\textbf{Schematic layout of the multi-scale quantum agrivoltaic digital twin platform.} Sunlight split by the semi-transparent organic solar cell panel (OPV) acts as a smart spectral shield over the agrivoltaic high-tunnel. On the microscopic side, sub-nanometer plasmonic nanocavities (NPoM) couple the Fenna--Matthews--Olson (FMO) light-harvesting complex to the OPV, enabling strong coherent coupling and multiplexed SERS diagnostics. Telemetry from the in situ sensor node is encrypted and transmitted via a BB84 fiber-optic quantum key distribution (QKD) link to a cloud server.}
\label{fig:system_setup}
\end{figure}

% =============================================================================
% REFERENCES
% =============================================================================

\section*{References}

\bibliographystyle{unsrt}
\bibliography{references}

\end{document}
